\begin{figure}[t]
\centering
\includegraphics[width=\textwidth]{fig2d_lie_algebra_symmetry.png}
\caption{\textbf{Dynamical Lie algebra analysis and effective dimension.} 
\textbf{a}, Commutator heatmap for the 10-qubit hardware-efficient ansatz. 
The full DLA has dimension $6.6\times10^5$; adaptive initialization restricts 
dynamics to a low-dimensional subalgebra (blue box). 
\textbf{b}, DLA dimension scaling with qubit count. Full SU($4^n$) DLA (gray dashed) 
grows as $4^n$. Hardware-efficient DLA (blue) saturates at $\sim2\times10^6$ for 
$n=100$ due to connectivity constraints. The accessible DLA (red), computed from 
experimental gradient ratios, is dramatically smaller—at $n=100$, $D_{\text{eff}} = 4.78\times10^{-6}$, 
representing a $10^{57}\times$ reduction in effective Hilbert space dimension. 
\textbf{c}, Schematic illustration. Adaptive initialization localizes the quantum 
dynamics to an exponentially small region of the hardware DLA, directly enabling 
the $4.18\times10^{11}\times$ gradient improvement observed experimentally. 
This demonstrates that barren plateaus are not fundamental—they are a consequence 
of initializing circuits to explore unnecessarily large regions of the DLA.}
\label{fig:fig2d}
\end{figure}